\scanpage{15}
\margin{$-9-$\qquad N$^\circ$ 234}
considéré comme espace vectoriel à droite sur $\varphi(A)$), et les $x_i\ne0$ appartiennent à $E$. Pour tout $y\in E$, on a alors la relation $\unsure{\Phi'|\Phi^\circ}(x',1\otimes y)=0$ qui s'écrit $\sum_{i=1}^n\alpha_i\varphi(\Phi(x_i,y))=0$, et entraîne donc $\Phi(x_i,y)=0$ pour $1\le i\le n$, d'où le corollaire.

4. \uline{Quelques identités}.

Soit $\Phi$ une forme sesquilinéaire sur $E$. Posons $\unsure{\Phi|\varphi}(x,x)=Q(x)$. On a
\[
\begin{aligned}
(9)\qquad Q(x+y)&=Q(x)+\Phi(x,y)+\Phi(y,x)+Q(y)\\
(10)\qquad Q(x-y)&=Q(x)-\Phi(x,y)-\Phi(y,x)+Q(y)
\end{aligned}
\]
d'où
\[
(11)\qquad 2(\Phi(x,y)+\Phi(y,x))=Q(x+y)-Q(x-y).
\]
Soit $\lambda\in A$ ; en remplaçant $y$ par $\lambda y$ dans (11), il vient
\[
(12)\qquad 2(\Phi(x,y)\lambda^J+\lambda\Phi(y,x))=Q(x+\lambda y)-Q(x-\lambda y)
\]
et on tire de (11) et (12)
\[
\begin{aligned}
(13)\qquad 2(\lambda\Phi(x,y)-\Phi(x,y)\lambda^J)
&=\lambda Q(x+y)-\lambda Q(x-y)-Q(x+\lambda y)+{}\\
&\qquad{}+Q(x-\lambda y).
\end{aligned}
\]
Supposons en particulier que $A$ soit une extension quadratique $K(i)$ d'un anneau commutatif $K$, avec $i^2=-1$ (chap. II, \S 7, n$^\circ$7), et que $J$ soit l'automorphisme $\alpha+\beta i\longrightarrow\alpha-\beta i$ de $A$ ; en faisant $\lambda=i$ dans (13), on obtient
\[
(14)\qquad 4\Phi(x,y)=Q(x+y)-Q(x-y)+iQ(x+iy)-iQ(x-iy).
\]
\margin{\uline{Remarque}}
On notera que les formules précédentes sont encore \deleted{applicables}\inserted{valables} si $\Phi$ est une application de $E\times E$ dans un $A$-bimodule $F$ (\uline{à gauche et à droite}), vérifiant les relations (1), (2), (3), (4).
\margin{?}
\margin{14}
